% French abstract (résumé) -- Mémoire de Master. Same content as abstract.tex, in French.
\thispagestyle{empty}
\vspace*{4cm}
\begin{center}{\Huge\bfseries Résumé}\end{center}
\vskip 1cm
\noindent Un essaim de drones commandé à la voix a besoin d'un analyseur qui transforme un énoncé
transcrit en une commande exécutable sur le matériel embarqué du véhicule, où chaque paramètre
supplémentaire coûte de la latence et où une sortie erronée est une erreur de pilotage. Ce mémoire
étudie le compromis entre précision et efficacité parmi de petits modèles de langage de 0,36 à
1,2~milliard de paramètres, affinés pour l'analyse de commandes d'essaims de drones et quantifiés
pour une exécution sur le seul processeur central d'un Raspberry~Pi~5.

\vspace{0.3cm}
\noindent Un schéma de commandes figé est associé à une grammaire qui lui correspond et contraint le
décodage, de sorte que la validité structurelle est une propriété du décodeur. Une chaîne de
constitution des données qui part des étiquettes produit les données d'entraînement et de test,
avec des partitions attribuées par famille de gabarits et un jeu de test de référence de
200~commandes enregistrées. Quatre modèles sont affinés selon une même recette~; trois sont
quantifiés à deux niveaux, leurs artefacts Q4\_K\_M sont chronométrés sur la carte, et le
quatrième, de taille appariée à l'un d'eux, mesure une différence entre familles de modèles à
taille fixée. Une règle de sélection fixée
avant l'obtention de tout résultat désigne la configuration à déployer.

\vspace{0.3cm}
\noindent Parmi les trois modèles chronométrés sur la carte, le compromis en Q4\_K\_M se réduit à
deux configurations non dominées, SmolLM2-360M et Qwen2.5-0.5B, et la règle déploie Qwen2.5-0.5B en
Q4\_K\_M, avec une correspondance exacte de 0,935 sur le jeu de référence, évaluée sur le texte de
référence des commandes. SmolLM2-360M est plus rapide et moins précis de 17,5~points de pourcentage, une
différence significative~; Llama-3.2-1B est plus lent, et sa précision inférieure sur le jeu de
référence n'est pas statistiquement établie. La sélection se fait sur la précision, face à des
contraintes redéfinies d'après les mesures~: la configuration retenue dépasse de 622~ms le budget
de latence de bout en bout de 2\,500~ms, et ce budget est redéfini à la valeur mesurée de
3\,122~ms. Aucune configuration ne satisfait les exigences de taux d'échec sûr et de taux de fausses
commandes, et ces deux écarts relèvent d'une même tendance~: lorsqu'une configuration se trompe ou
que l'entrée est hors domaine, elle émet une commande bien formée au lieu de s'abstenir. Un seuil de
confiance appliqué à la sortie du décodeur est la première étape proposée.

\vspace{1cm}
\providecommand{\motscles}[1]{\textbf{\textit{Mots-clés---}} #1}
\motscles{petits modèles de langage~; inférence embarquée~; quantification~; décodage contraint
par grammaire~; analyse de commandes~; Raspberry Pi~5~; essaim de drones~; front de Pareto.}
\newpage
